\section{Otentikasi dan Non-Repudiation}
Dalam keamanan informasi, mengetahui \textit{siapa} yang mengirimkan pesan sama pentingnya dengan menjaga kerahasiaan isi pesan tersebut.

\subsection{Otentikasi}
Otentikasi adalah proses pembuktian identitas seorang pengguna atau sistem. Terdapat tiga faktor utama otentikasi:
\begin{enumerate}
    \item \textbf{Something you know}: Kata sandi (\textit{password}), PIN.
    \item \textbf{Something you have}: Token fisik, kartu pintar (\textit{smart card}), sertifikat digital.
    \item \textbf{Something you are}: Biometrik (sidik jari, scan retina).
\end{enumerate}

\subsection{Non-Repudiation (Nir-Sangkalan)}
\textit{Non-repudiation} adalah jaminan bahwa pengirim pesan tidak dapat menyangkal telah mengirimkan pesan tersebut, dan penerima tidak dapat menyangkal telah menerima pesan tersebut. Hal ini dicapai melalui penggunaan tanda tangan digital.

\subsection{Otentikasi Entitas vs Otentikasi Pesan}
\textbf{Otentikasi entitas} membuktikan identitas pihak yang berkomunikasi pada saat sesi dibangun (misalnya \textit{login}), umumnya dengan tantangan \textit{nonce} agar tahan \textit{replay}. \textbf{Otentikasi pesan} membuktikan bahwa tiap pesan yang diterima benar-benar berasal dari entitas tersebut dan belum diubah; ini tugas MAC atau tanda tangan digital \citep{menezes1996}. Keduanya perlu: sesi yang sudah terotentikasi tetap dapat menerima pesan palsu bila tidak ada otentikasi per pesan.

\textit{Non-repudiation} menuntut bukti yang dapat diverifikasi pihak ketiga, sehingga hanya dapat dicapai dengan tanda tangan digital (dan infrastruktur kunci publik), bukan dengan MAC yang kuncinya simetris \citep{stallings2017}.

